% Zumdahl Chapter 6 free-response set -- 2 long (10) + 3 short (4) = 32 pts.
\documentclass{apchem-exam}

\apcunitword{Chapter}
\apcunit{6}
\apcunittitle{Thermochemistry}
\apcdoctitle{Free-Response Set}
\apcsubtitle{Zumdahl \S6.1--\S6.4 \textbullet\ 5 questions \textbullet\ 32 points \textbullet\ 50 minutes}

\begin{document}
\apctitleblock

\begin{apcnote}[Scope]
Covers \textbf{CED 6.1--6.4} (\S6.1--6.2), \textbf{6.6} (\S6.2),
\textbf{6.8} (\S6.4) and \textbf{6.9} (\S6.3). Two neighbouring CED
topics are sourced from other chapters and are not examined here:
\textbf{6.5} (energy of phase changes) comes from \S10.8, and
\textbf{6.7} (bond enthalpies) from \S8.8. \S6.5--6.6 of Zumdahl (energy
sources) are context and are not assessed.

Take the specific heat of water as
\SI{4.18}{\joule\per\gram\per\celsius} and assume dilute aqueous
solutions have the density and specific heat of water.
\end{apcnote}

\examsection{II}{Free Response}{5 questions \textbullet\ 32 points \textbullet\ 50 minutes}{\calculatorok}

% =====================================================================
\frq{long}{10}{Zumdahl \S6.2 \textbullet\ CED 6.4, 6.6 \textbullet\ calorimetry of a neutralization}

A student mixes \SI{50.0}{\milli\litre} of \SI{1.00}{\Molar} \ce{HCl}
with \SI{50.0}{\milli\litre} of \SI{1.00}{\Molar} \ce{NaOH} in a
polystyrene cup. Both solutions start at \SI{21.0}{\celsius}. The highest
temperature reached is \SI{27.8}{\celsius}.

\begin{parts}
  \item Write the net ionic equation for the reaction.
        \answerlines[2]{\ce{H+(aq) + OH^-(aq) -> H2O(l)}}
  \item Calculate the heat released by the reaction. \workspace[2.8cm]
        \keyonly{\begin{apcanswer}Total mass of solution
        $= 50.0 + 50.0 = \SI{100.0}{\gram}$ (assuming the density of
        water). $\Delta T = 27.8 - 21.0 = \SI{6.8}{\celsius}$.

        $q = mc\,\Delta T = (100.0)(4.18)(6.8) = \mathbf{2.8\times10^{3}}$
        J $= \SI{2.84}{\kilo\joule}$ absorbed by the solution, so
        \SI{2.84}{\kilo\joule} was released by the
        reaction.\end{apcanswer}}
  \item Calculate the enthalpy of neutralization in
        \si{\kilo\joule\per\mole}, and state its sign.
        \workspace[2.8cm]
        \keyonly{\begin{apcanswer}$n(\ce{H+}) = 0.0500 \times 1.00 =
        \num{0.0500}$~mol, and \ce{OH-} is present in exactly the same
        amount, so \num{0.0500}~mol of water forms.

        $\Delta H = -\dfrac{2.84}{0.0500} =
        \mathbf{-56.8}$~kJ/mol

        The sign is \textbf{negative}: the temperature rose, so the
        reaction released energy to the solution and is
        exothermic.\end{apcanswer}}
  \item The student repeats the experiment using
        \SI{100.0}{\milli\litre} of each solution at the same
        concentration. State what happens to the temperature change and to
        $\Delta H$ per mole, and explain.
        \answerlines[3]{The temperature change is essentially
        \textbf{unchanged} (about \SI{6.8}{\celsius}): twice as much heat
        is released, but it warms twice as much solution. $\Delta H$ per
        mole is also \textbf{unchanged} --- it is an intensive quantity
        that describes the reaction itself, not how much of it was run.}
  \item Give one reason why the measured magnitude of $\Delta H$ is
        typically smaller than the accepted value.
        \answerlines[2]{Heat is lost to the surroundings --- to the cup,
        the thermometer and the air --- so the measured temperature rise,
        and hence the calculated heat, is too small. (Accept also: the
        heat capacity of the calorimeter itself is ignored.)}
\end{parts}

\begin{rubric}
  \rpoint{1}{(a) \ce{H+ + OH^- -> H2O}}
  \rpoint{3}{(b) 1 pt total mass \SI{100.0}{\gram}; 1 pt
    $\Delta T = \SI{6.8}{\celsius}$; 1 pt \SI{2.8e3}{\joule}. Using
    \SI{50.0}{\gram} caps this part at 2}
  \rpoint{3}{(c) 1 pt \num{0.0500}~mol; 1 pt dividing heat by moles; 1 pt
    the negative sign with an exothermic justification. A positive
    $\Delta H$ with a rising temperature earns 0 for the sign point}
  \rpoint{2}{(d) 1 pt $\Delta T$ unchanged with the reasoning; 1 pt
    $\Delta H$ per mole unchanged because it is intensive}
  \rpoint{1}{(e) heat loss to the surroundings, or calorimeter heat
    capacity neglected}
\end{rubric}

% =====================================================================
\frq{long}{10}{Zumdahl \S6.3--\S6.4 \textbullet\ CED 6.8, 6.9 \textbullet\ Hess's law and enthalpies of formation}

Consider the following data at \SI{298}{\kelvin}:
\begin{center}
\begin{tabular}{lr}
\hline
\textbf{Reaction} & $\Delta H^\circ$ (kJ) \\
\hline
\ce{C(s) + 1/2O2(g) -> CO(g)}   & $-110.5$ \\
\ce{CO(g) + 1/2O2(g) -> CO2(g)} & $-283.0$ \\
\hline
\end{tabular}
\end{center}

\begin{parts}
  \item Use Hess's law to determine $\Delta H^\circ$ for
        \ce{C(s) + O2(g) -> CO2(g)}. \workspace[2.8cm]
        \keyonly{\begin{apcanswer}The two reactions add directly: the
        \ce{CO} produced by the first is consumed by the second and
        cancels.

        $\Delta H^\circ = (-110.5) + (-283.0) =
        \mathbf{-393.5}$~kJ\end{apcanswer}}
  \item Explain why this value is also the standard enthalpy of formation
        of \ce{CO2}. \answerlines[3]{The reaction forms one mole of
        \ce{CO2} from its elements, each in its standard state ---
        graphite and \ce{O2} gas --- at \SI{298}{\kelvin} and
        \SI{1}{\atm}. That is precisely the definition of a standard
        enthalpy of formation.}
  \item State the value of $\Delta H_f^\circ$ for \ce{O2(g)}, with a
        reason. \answerlines[2]{\textbf{Zero}. \ce{O2} gas is the standard
        state of the element oxygen, and forming an element from itself
        involves no change.}
  \item Explain why Hess's law works --- that is, why enthalpy changes may
        be added in this way. \workspace[2.6cm]
        \keyonly{\begin{apcanswer}Enthalpy is a \textbf{state function}:
        its value depends only on the current state of the system, not on
        the path taken to reach it. The enthalpy difference between
        reactants and products is therefore fixed, so any sequence of
        steps connecting them must sum to the same total. This is why an
        enthalpy change that is hard to measure directly can be obtained
        from steps that are easy to measure.\end{apcanswer}}
  \item If the second reaction were reversed, state the value of its
        $\Delta H^\circ$ and explain. \answerlines[2]{$+283.0$~kJ.
        Reversing a reaction reverses the direction of energy flow, so the
        magnitude is unchanged and the sign flips.}
\end{parts}

\begin{rubric}
  \rpoint{2}{(a) 1 pt recognizing the equations add with \ce{CO}
    cancelling; 1 pt $-393.5$~kJ}
  \rpoint{2}{(b) 1 pt one mole of product formed; 1 pt from elements in
    their standard states}
  \rpoint{2}{(c) 1 pt zero; 1 pt because \ce{O2} is the element's standard
    state}
  \rpoint{2}{(d) 1 pt enthalpy is a state function; 1 pt therefore
    path-independent, so steps sum. ``Because energy is conserved'' alone
    earns 1}
  \rpoint{2}{(e) 1 pt $+283.0$~kJ; 1 pt sign reverses on reversing the
    reaction}
\end{rubric}

% =====================================================================
\frq{short}{4}{Zumdahl \S6.2 \textbullet\ CED 6.4 \textbullet\ specific heat}

A \SI{125}{\gram} block of an unknown metal at \SI{95.0}{\celsius} is
dropped into \SI{200.0}{\gram} of water at \SI{22.0}{\celsius}. The final
temperature of both is \SI{27.1}{\celsius}.

\begin{parts}
  \item Calculate the specific heat of the metal. \workspace[3.2cm]
        \keyonly{\begin{apcanswer}Heat gained by the water:

        $q = (200.0)(4.18)(27.1 - 22.0) = \SI{4264}{\joule}$

        The metal lost the same amount, cooling from \SI{95.0}{\celsius}
        to \SI{27.1}{\celsius}, a fall of \SI{67.9}{\celsius}:

        $c = \dfrac{4264}{(125)(67.9)} =
        \mathbf{0.502}$~J/g\,\si{\celsius}\end{apcanswer}}
  \item State the assumption about the surroundings that this calculation
        requires. \answerlines[2]{That no heat is exchanged with anything
        else --- all the heat lost by the metal is gained by the water,
        with none lost to the container or the air.}
\end{parts}

\begin{rubric}
  \rpoint{3}{(a) 1 pt heat gained by water; 1 pt equating it to heat lost
    by the metal with the metal's own $\Delta T$ of \SI{67.9}{\celsius};
    1 pt \SI{0.502}{}~J/g\,\si{\celsius}. Using the water's $\Delta T$ for
    the metal caps this at 1}
  \rpoint{1}{(b) a statement that the calorimeter is isolated / no heat
    lost to the surroundings}
\end{rubric}

% =====================================================================
\frq{short}{4}{Zumdahl \S6.1 \textbullet\ CED 6.1--6.3 \textbullet\ energy diagrams and heat flow}

\begin{parts}
  \item When ammonium nitrate dissolves in water, the beaker becomes cold
        to the touch. State the sign of $\Delta H$ for the dissolution and
        explain the direction of heat flow.
        \answerlines[3]{$\Delta H$ is \textbf{positive} --- the process is
        endothermic. Energy flows \emph{from} the water \emph{into} the
        dissolving solid, so the water loses energy and its temperature
        falls, which is what the hand detects as cold.}
  \item On an energy diagram for this process, state whether the products
        lie above or below the reactants, and why.
        \answerlines[2]{\textbf{Above.} The system absorbed energy, so the
        products are at higher enthalpy than the reactants; the vertical
        gap between them is $\Delta H$.}
\end{parts}

\begin{rubric}
  \rpoint{2}{(a) 1 pt positive / endothermic; 1 pt heat flows from the
    water into the system. Saying ``the beaker gives off cold'' earns 0}
  \rpoint{2}{(b) 1 pt products above reactants; 1 pt because energy was
    absorbed}
\end{rubric}

% =====================================================================
\frq{short}{4}{Zumdahl \S6.4 \textbullet\ CED 6.8 \textbullet\ using enthalpies of formation}

For the combustion of methane,
\ce{CH4(g) + 2O2(g) -> CO2(g) + 2H2O(l)}, use
$\Delta H_f^\circ$: \ce{CH4(g)} $= -74.8$, \ce{CO2(g)} $= -393.5$,
\ce{H2O(l)} $= -285.8$ \si{\kilo\joule\per\mole}.

\begin{parts}
  \item Calculate $\Delta H^\circ$ for the reaction. \workspace[3cm]
        \keyonly{\begin{apcanswer}$\Delta H^\circ = \sum \Delta H_f^\circ
        (\text{products}) - \sum \Delta H_f^\circ (\text{reactants})$

        Products: $(-393.5) + 2(-285.8) = \num{-965.1}$

        Reactants: $(-74.8) + 2(0) = \num{-74.8}$ \quad
        (\ce{O2} is an element in its standard state)

        $\Delta H^\circ = -965.1 - (-74.8) =
        \mathbf{-890.3}$~kJ\end{apcanswer}}
\end{parts}

\begin{rubric}
  \rpoint{4}{2 pts the products-minus-reactants relationship applied with
    the coefficient 2 on water; 1 pt using zero for \ce{O2}; 1 pt
    $-890.3$~kJ. Reversing the subtraction gives $+890.3$ and earns at
    most 2}
\end{rubric}

\examtotal

\end{document}
